\documentclass[11pt]{article}
\usepackage{amsmath,amssymb,geometry,hyperref}
\geometry{margin=1in}
\newcommand{\Tr}{\operatorname{Tr}}
\begin{document}
\title{An analytical optimal second-order AC subspace of eighteen qubits}
\author{Catalogue calculation}\date{30 October 2026}\maketitle
Use normalized Dicke states $D_r=|D_{18}^{(r)}\rangle$ with
$J_zD_r=(9-r)D_r$. We construct a four-dimensional subspace for which every
normalized complex vector has two-qubit reduction $I_3/3$.

\section{Explicit radical basis}
\[
\begin{aligned}|u_1\rangle={}&a(D_0+D_{18})+b(D_3+D_{15})\\&+c(D_6+D_{12})+dD_9,\end{aligned}
\]
\[
\begin{aligned}a&=\frac{\sqrt{905}}{81}\\b&=- \frac{7 \sqrt{507705}}{49232} - \frac{9925 \sqrt{9231} i}{3987792}\\c&=\frac{1081 \sqrt{4200105}}{90722268} + \frac{145 \sqrt{9240231} i}{1120028}\\d&=\frac{99 \sqrt{40001}}{160004} - \frac{4481 \sqrt{2200055} i}{12960324}\end{aligned}
\]
\[
|u_2\rangle=e(D_0-D_{18})+f(D_3-D_{15})+g(D_6-D_{12}),
\]
\[
\begin{aligned}e&=\frac{1}{3}\\f&=\frac{5 \sqrt{561}}{816} + \frac{41 \sqrt{255} i}{2448}\\g&=\frac{\sqrt{4641}}{13923} - \frac{5 \sqrt{255255} i}{4641}\end{aligned}
\]
\[
|u_3\rangle=\sum_{j=0}^5h_jD_{1+3j},
\]
\[
\begin{aligned}h_0&=\frac{4 \sqrt{5}}{27}\\h_1&=- \frac{7 \sqrt{374}}{816} + \frac{263 \sqrt{170} i}{7344}\\h_2&=- \frac{1249 \sqrt{2210}}{190944} + \frac{95 \sqrt{4862} i}{21216}\\h_3&=- \frac{11 \sqrt{1105}}{1768} + \frac{175 \sqrt{2431} i}{47736}\\h_4&=- \frac{473 \sqrt{595}}{51408} - \frac{25 \sqrt{1309} i}{5712}\\h_5&=- \frac{\sqrt{1870}}{102} - \frac{55 \sqrt{34} i}{918}\end{aligned}
\]
\[
|u_4\rangle=\sum_{j=0}^5h_jD_{17-3j}.
\]

Every unlisted coefficient is zero. The first vector is reflection-even,
the second reflection-odd, and the fourth is the reversed third vector.
These support and reflection conventions are chosen for a compact basis;
they are not assumptions in the upper-bound theorem.

\section{Exact whole-subspace verification}
For any coefficient vector $c$, define its Dicke splitting matrix
\[
 (M_t(c))_{ab}=c_{a+b}
 \sqrt{\frac{\binom ta\binom{18-t}b}{\binom{18}{a+b}}},
 \qquad 0\le a\le t,\quad0\le b\le18-t.
\]
The ordered reduction of $|u_i\rangle\langle u_j|$ is
$M_t(u_i)M_t(u_j)^\dagger$. Substitution of the displayed radical
coefficients gives, with exact arithmetic,
\[
 \langle u_i,u_j\rangle=\delta_{ij},\qquad
 M_2(u_i)M_2(u_j)^\dagger=\frac{\delta_{ij}}3 I_3
 \quad(1\le i,j\le4).
\]
The bundled verifier reconstructs all sixteen ordered cross matrices and
checks all 144 complex entries, including every off-diagonal basis pair.
These explicit radical identities prove the claim independently of the
numerical search that suggested the coefficients.
Consequently, for arbitrary complex $\alpha_i$ satisfying
$\sum_i|\alpha_i|^2=1$, the vector
$|\psi\rangle=\sum_i\alpha_i|u_i\rangle$ has
$\rho_2=I_3/3$.

A useful way to check the diagonal reductions is to note the exact
probability lists, in ascending order on the respective supports:
\[
\begin{aligned}
 u_1:&\quad (1810,885,2040,3652,2040,885,1810)/13122,\\
 u_2:&\quad (6,5,16,16,5,6)/54,\\
 u_3:&\quad (80,179,140,55,55,220)/729,\\
 u_4:&\quad (220,55,55,140,179,80)/729.
\end{aligned}
\]
Each has the required first and second spin moments, and support spacing
three eliminates its first- and second-step self coherences. The cross
conditions require the displayed phases as well; diagonal moments alone
would not certify the whole subspace.

The third-qubit reduced purities are
\[
 \Tr\rho_3(u_1)^2=\frac{39467}{147968},\quad
 \Tr\rho_3(u_2)^2=\frac{37267}{147968},\quad
 \Tr\rho_3(u_3)^2=\Tr\rho_3(u_4)^2=\frac{86707}{332928}.
\]
Each exceeds $1/4$, so every displayed basis vector has exactly AC order
two. This does not classify the orders of all their superpositions.

\section{Why the dimension is maximal}
Let $P$ project onto any second-order AC subspace of spin $j=9$, with
dimension $k$. Its polarized moment identities give
\[
 PJ_aP=0,\qquad P\frac{J_aJ_b+J_bJ_a}{2}P=30\delta_{ab}P.
\]
Couple to an auxiliary spin $s=3/2$, and let $X=\mathbf J\cdot\mathbf S$.
The compression of $X$ to $\mathcal C\otimes\mathbb C^4$ is zero, while
that of $X^2$ is $(225/2)I$. Hence the Hermitian operator
\[
 W=X^2-7X-\frac{225}{2}I
\]
has zero compression to this $4k$-dimensional subspace.
On total spins $L=15/2,17/2,19/2,21/2$, its eigenvalues are respectively
\[
 \frac{435}{2},\quad-\frac{99}{4},\quad-\frac{249}{2},\quad-\frac{99}{4}.
\]
There are sixteen positive and sixty negative eigenvalues, and no zeros.
Projection of a totally isotropic subspace into the positive eigenspace
is injective: a vector with zero positive component would have strictly
negative quadratic form unless it vanished. Therefore
\[
 4k\le16,\qquad k\le4.
\]
The radical basis attains four, proving
\[
 \boxed{K(18,2)=4.}
\]
This is the specialization of the general auxiliary-spin bound
$4k\le N-2$, and applies to arbitrary complex subspaces.

\section{Reconstruction and scope}
The starting numerical example is John Martin's
\texttt{\(\Pi\)\_18\_4\_2} in the supplied \texttt{ACLiege.jl}, line 966.
The source span nearly commutes with the residue-three decomposition.
Its components have dimensions two, one and one in Dicke residues zero,
one and two. A common rotation aligns a half-turn with coefficient reversal;
a unitary basis change then selects the even and odd vectors in residue
zero and the reflected pair in residues one and two.

The rational probability lists above allow refinement of the ten remaining
phases. Their real and imaginary parts reconstruct as square roots of
rational numbers. The displayed radicals are certified by the exact
whole-subspace identities, without interpreting decimal phase matching
as a proof.

For the coefficient multiplier $c_r\mapsto e^{-ir\beta/2}c_r$, with
$\beta\simeq0.734100176590464$, the original numerical subspace projector
matches the analytical projector within $8.62\times10^{-8}$ in Frobenius
norm. This is approximate association to the numerical source; literal
componentwise equality to its basis is not asserted. No classification of
all optimal subspaces or claim of literature priority is made.
\end{document}
